\section{Discussion}
The principal empirical result is lower rolling absolute coverage error for the two specified controllers on most of the selected streams. It is narrower than either a new conformal theorem or a general claim about forecasting architecture. The interval rules respond to accumulated past misses, while their shared point model remains unchanged. This mechanism makes a coverage improvement plausible, but the proper-score comparison shows why plausibility should not be mistaken for complete uncertainty improvement. Slightly lower aggregate interval score coexists with worse scores on a majority of cases.

The difference between the two coverage summaries also changes the interpretation of magnitude. The previously emphasized sevenfold ratio uses a statistic that averages signed rolling deviations before applying the absolute value. Its improvement can reflect balancing undercoverage and overcoverage at different times. Averaging absolute rolling deviations instead yields reductions near one fifth on these streams. Neither number is erroneous when correctly defined; the error would be to use the smaller statistic's ratio as if it described uniformly accurate local coverage. The deterministic inequality and rolling-boundary weights explain that distinction without relying on a significance test.

The sensitivity grid identifies dependence on the base and nominal level. It is useful because a controller tuned around one error distribution can behave differently around another. It cannot establish a generally optimal gain, especially when the original gain rationale is unavailable and the same streams are used for every cell. A stronger follow-up would prospectively fix authentic-data collections, identify a separate tuning budget, and include faithful published adaptive and score-predictive comparators. AgACI, strongly adaptive methods, decaying-step procedures and SPCI address relevant design questions, but none is numerically evaluated by the historical controller labels.

For research and teaching, the practical value of the replay is its inspectable information order and complete numerical state. A reader can locate the forecast made before a target arrived, reconstruct its buffer, separate its point error from interval changes, and check the chosen coverage statistic. For an operational forecasting service, additional questions remain: the cost of native updates, delayed or revised targets, missing observations, covariate arrival, simultaneous horizons and a loss attached to decisions. The current one-step full-prefix replay neither answers nor rules out those extensions. Its bounded findings provide a reproducible reference from which those better-specified experiments can proceed.

